\section{The auction stage: who pays what}\label{sec:03_auction}

This section derives equations \paper{(6)}--\paper{(8)}, Proposition~1, Proposition~1(c) and the surplus
split. Everything here holds \emph{before} trading and entry. It is pure auction algebra, and it does not
depend on the order slot.

\subsection{Rules of a second-price auction with reserve}\label{sec:03_auction-rules}

\begin{itemize}
\item A bid is \emph{admissible} if it is at least \(p\). A bid equal to \(p\) is admissible.
\item The highest admissible bidder wins. There is no sale without an admissible bid.
\item The winner pays \(\max\{p,\text{highest competing admissible bid}\}\).
\item Truthful bids are weakly dominant. So each prepared bidder whose value meets the reserve bids its
value: the incumbent bids \(R\) if \(R\ge p\), and a prepared challenger bids \(\theta\) if \(\theta\ge p\).
A bidder whose value is below \(p\) does not bid.
\end{itemize}

\paragraph{Why truthful bids are weakly dominant} \paper{App.~A.1}. Fix a bidder with value \(v\). Let
\(Y\ge0\) be the highest competing admissible bid, with \(Y=0\) if there is none. Any admissible bid that
wins pays \(\max\{p,Y\}\). So a bid yields either \(v-\max\{p,Y\}\) (a win) or zero (a loss).
\begin{enumerate}
\item Let \(v\ge p\), and consider the bid \(v\).
\begin{itemize}
\item If \(v>\max\{p,Y\}\), the bid wins and earns \(v-\max\{p,Y\}>0\).
\item If \(v<\max\{p,Y\}\), then \(Y>v\), so the bid loses and earns zero.
\item If \(v=\max\{p,Y\}\), a win pays \(v\), so the bid earns zero.
\end{itemize}
In each case the bid \(v\) earns \((v-\max\{p,Y\})_+\), the larger of the two possible payoffs.
\item Let \(v<p\). Every bid that wins pays at least \(p>v\). So no bid earns more than zero, which is the
payoff of no bid.
\end{enumerate}
The argument holds at every \(Y\). So it needs no belief about the rival. None of the payoffs below
depends on a conjecture about bid shading. This verifies condition (E6) for every tuple.

\subsection{Case-by-case outcomes}\label{sec:03_auction-cases}

Fix the incumbent's value \(R\). Table~\ref{tab:03_auction-cases} lists the winner and the payment in
each case of the benchmark. Value ties have probability zero, because \(F\) is continuous.

\begin{table}[!htbp]
\centering\small
\begin{tabular}{@{}llll@{}}
\toprule
Challenger & Region of \(R\) & Winner & Target receives\\
\midrule
None & \(R<p\) & no sale & \(0\)\\
None & \(R\ge p\) & incumbent & \(p\)\\
\midrule
High (\(\theta=h\)) & \(R<p\) & challenger & \(p\)\\
High & \(p\le R\le r\) & challenger & \(R\)\\
\midrule
Low (\(\theta=\ell\)) & \(R<p\) & challenger & \(p\)\\
Low & \(p\le R<\ell\) & challenger & \(R\)\\
Low & \(\ell<R\le r\) & incumbent & \(\ell\)\\
\bottomrule
\end{tabular}
\caption{Realized sale outcomes in the benchmark. A high-value challenger always wins because
\(h>r\ge R\).}
\label{tab:03_auction-cases}
\end{table}

Compact forms of what the target receives:
\[
\text{no entry: } p\,\1\{R\ge p\},\qquad
\text{high challenger: } \max\{p,R\},\qquad
\text{low challenger: } \max\{p,\min(R,\ell)\}.
\]
For a general value \(\theta\) \paper{Sec.~3.1}:
\begin{itemize}
\item if \(\theta\ge p\), the target receives \(\max\{p,\min(R,\theta)\}\), and the challenger keeps
\((\theta-\max\{p,R\})_+\);
\item if \(\theta<p\), the challenger cannot bid. The target receives \(p\,\1\{R\ge p\}\), as without entry,
and the challenger keeps nothing.
\end{itemize}

\begin{intuition}[the one line that drives the paper]
If \(R\le\ell\), both challenger values give the target the same payment,
\(\max\{p,R\}\). If \(R>\ell\), a high-value challenger pays \(R\), but a low-value challenger loses, and
the incumbent pays \(\ell\). So the gap in target proceeds is
\[
(\text{proceeds with }h)-(\text{proceeds with }\ell)=(R-\ell)_+ .
\]
A stronger incumbent puts more mass on large \(R\). This raises the gap. The same shift raises what a
challenger must pay to win. This lowers its profit.
\end{intuition}

\begin{figure}[htbp]
\centering
\begin{tikzpicture}[x=3.1cm, y=1.25cm, font=\small]
\fill[intuit!12] (1,1) -- (3,3) -- (3,1) -- cycle;
\draw[-{Stealth[length=2mm]}] (0,0) -- (3.25,0) node[right] {\(R\)};
\draw[-{Stealth[length=2mm]}] (0,0) -- (0,3.35) node[above] {target receives};
\node[below] at (1.75,-0.32) {incumbent value \(R\)};
\draw[gray, densely dotted, very thick] (0,0) -- (0.5,0);
\draw[gray, densely dotted, very thick] (0.5,0.5) -- (3,0.5);
\draw[intuit, very thick] (0,0.5) -- (0.5,0.5) -- (3,3);
\draw[panel, very thick, dashed] (0,0.5) -- (0.5,0.5) -- (1,1) -- (3,1);
\foreach \x/\l in {0.5/p, 1/\ell, 3/r} \draw (\x,0.05) -- (\x,-0.05) node[below] {\(\l\)};
\node[below] at (0,-0.05) {\(0\)};
\foreach \y/\l in {0.5/p, 1/\ell} \draw (0.04,\y) -- (-0.04,\y) node[left] {\(\l\)};
\node[font=\footnotesize, intuit, anchor=west] at (0.75,2.2) {high-value challenger};
\node[font=\footnotesize, panel] at (2.25,1.17) {low-value challenger};
\node[font=\footnotesize, gray] at (2.25,0.34) {no entry};
\node[font=\footnotesize] at (2.6,1.75) {\((R-\ell)_+\)};
\end{tikzpicture}
\caption{Target proceeds as a function of \(R\) (\(p=1/2\), \(\ell=1\), \(r=3\)). The two challenger lines
coincide for \(R\le\ell\). The shaded gap is \((R-\ell)_+\); its expectation is \(\DT\). The dotted line
is the no-entry payment \(p\,\1\{R\ge p\}\); both challenger lines lie on or above it.}
\label{fig:03_auction-proceeds}
\end{figure}

\subsection{Derivation of the six payoff terms}\label{sec:03_auction-terms}

\(R\) is uniform on \([0,r]\), so its density is \(1/r\). We use three facts:
\[
\int_a^b c\,dy=c(b-a),\qquad \int_a^b y\,dy=\frac{b^2-a^2}{2},\qquad
\int_a^b(c-y)\,dy=\frac{(c-a)^2}{2}-\frac{(c-b)^2}{2}.
\]

\paragraph{(a) No entry, \(t_0\).}
\[
t_0=\E\bigl[p\,\1\{R\ge p\}\bigr]=p\Pr(R\ge p)=p\cdot\frac{r-p}{r}=p\Bigl(1-\frac pr\Bigr).
\]

\paragraph{(b) High challenger, \(t_H\).} The challenger always wins and pays \(\max\{p,R\}\). Split at
\(R=p\):
\[
t_H=\frac1r\Bigl[\int_0^p p\,dy+\int_p^r y\,dy\Bigr]=\frac1r\Bigl[p^2+\frac{r^2-p^2}{2}\Bigr]
=\frac1r\cdot\frac{r^2+p^2}{2}=\frac r2+\frac{p^2}{2r}.
\]

\paragraph{(c) High challenger's profit, \(g_H\).} It wins at every \(R\) and keeps \(h-\max\{p,R\}\). So
\[
g_H=\frac1r\int_0^r\bigl(h-\max\{p,y\}\bigr)\,dy=h-t_H=h-\frac r2-\frac{p^2}{2r}.
\]

\paragraph{(d) Low challenger, \(t_L\).} Split at \(R=p\) and \(R=\ell\):
\begin{align*}
t_L&=\frac1r\Bigl[\int_0^p p\,dy+\int_p^\ell y\,dy+\int_\ell^r\ell\,dy\Bigr]
=\frac1r\Bigl[p^2+\frac{\ell^2-p^2}{2}+\ell(r-\ell)\Bigr]\\
&=\frac1r\Bigl[\ell r-\frac{\ell^2}{2}+\frac{p^2}{2}\Bigr]=\ell-\frac{\ell^2-p^2}{2r}.
\end{align*}
The second line collects terms: \(p^2-p^2/2=p^2/2\), and \(\ell^2/2-\ell^2=-\ell^2/2\).

\paragraph{(e) Low challenger's profit, \(g_L\).} It wins only if \(R<\ell\), and keeps
\(\ell-\max\{p,R\}\):
\begin{align*}
g_L&=\frac1r\Bigl[\int_0^p(\ell-p)\,dy+\int_p^\ell(\ell-y)\,dy\Bigr]
=\frac1r\Bigl[(\ell-p)p+\frac{(\ell-p)^2}{2}\Bigr]\\
&=\frac{(\ell-p)(2p+\ell-p)}{2r}=\frac{(\ell-p)(\ell+p)}{2r}=\frac{\ell^2-p^2}{2r}.
\end{align*}
The second integral uses the third fact with \(c=\ell\), \(a=p\), \(b=\ell\).

\paragraph{(f) The spread, \(\DT\).} Subtract (d) from (b):
\[
\DT=t_H-t_L=\frac r2+\frac{p^2}{2r}-\ell+\frac{\ell^2-p^2}{2r}=\frac r2-\ell+\frac{\ell^2}{2r}
=\frac{r^2-2\ell r+\ell^2}{2r}=\frac{(r-\ell)^2}{2r}.
\]
The reserve \(p\) cancels. This is no accident: below \(\ell\) both values pay the same, so \(p\) moves
both terms equally. A direct check: \(\DT=\E[(R-\ell)_+]=\frac1r\int_\ell^r(y-\ell)\,dy=(r-\ell)^2/(2r)\).

\begin{resultbox}[equation (6) of the paper]
\[
t_0=p\Bigl(1-\frac pr\Bigr),\qquad t_H=\frac r2+\frac{p^2}{2r},\qquad t_L=\ell-\frac{\ell^2-p^2}{2r},
\]
\[
g_H=h-\frac r2-\frac{p^2}{2r},\qquad g_L=\frac{\ell^2-p^2}{2r},\qquad \DT(r)=\frac{(r-\ell)^2}{2r}.
\]
\end{resultbox}

\begin{panelnote}[three quick sanity checks]
\begin{itemize}
\item \(t_H+g_H=h\) and \(t_L+g_L=\ell\). Proof for \(\ell\): if \(R<\ell\), the challenger wins, and the
payment plus its profit equals \(\ell\). If \(R>\ell\), the target gets \(\ell\) and the challenger gets
\(0\). The sum is \(\ell\) in both cases. The incumbent's profit \(R-\ell\) is outside this sum.
\item \(w_L=t_L-t_0>0\), because \(\max\{p,\min(R,\ell)\}\ge p\ge p\,\1\{R\ge p\}\), and the second
inequality is strict on \(R<p\), an event of probability \(p/r>0\). So entry by either class raises
proceeds. At \(r_1=3\), \(w_L=11/24\) and \(w_H=9/8\).
\end{itemize}
\end{panelnote}

\subsection{Derivatives: equation (7)}\label{sec:03_auction-derivatives}

\paragraph{Spread.} Use the quotient rule on \(\DT=(r-\ell)^2/(2r)\):
\[
\DT'(r)=\frac{2(r-\ell)\cdot2r-(r-\ell)^2\cdot2}{4r^2}=\frac{(r-\ell)\,[2r-(r-\ell)]}{2r^2}
=\frac{(r-\ell)(r+\ell)}{2r^2}=\frac12-\frac{\ell^2}{2r^2}>0,
\]
since \(r>\ell\).

\paragraph{Profits.}
\[
g_H'(r)=-\frac12+\frac{p^2}{2r^2}<0\quad(\text{since }r>p),\qquad
g_L'(r)=-\frac{\ell^2-p^2}{2r^2}<0\quad(\text{since }\ell>p). \tag*{\paper{(7)}}
\]

\paragraph{Gross profit at a fixed belief.} \(B_r(\mu)=(1-\mu)g_L+\mu g_H\) is a weighted average with
fixed weights. So
\[
\frac{\partial B_r(\mu)}{\partial r}=(1-\mu)g_L'(r)+\mu g_H'(r)<0\qquad\text{for every }\mu\in[0,1].
\]
Also \(B_r\) rises in \(\mu\), because
\[
g_H-g_L=h-\frac r2-\frac{p^2}{2r}-\frac{\ell^2-p^2}{2r}=h-\frac r2-\frac{\ell^2}{2r}
>r-\frac r2-\frac{\ell^2}{2r}=\frac{r^2-\ell^2}{2r}>0\qquad(\text{use }h>r>\ell).
\]
At \(r_1=3\): \(\DT'=4/9\), \(g_H'=-35/72\), \(g_L'=-1/24\), and \(\partial B_r(1/2)/\partial r=-19/72\).

\paragraph{The preparation threshold.} With \(g_H>g_L\), solve \(B_r(\mu)=c_H\):
\begin{align*}
g_L+\mu(g_H-g_L)=c_H
&\iff \mu=\frac{c_H-g_L}{g_H-g_L}=:\tau(r),\\
c_H-g_L&=\frac{2rc_H-\ell^2+p^2}{2r},\qquad
g_H-g_L=\frac{2rh-r^2-\ell^2}{2r}.
\end{align*}
Cancel the common denominator \(2r\):
\begin{equation}\label{eq:03-tau}
\tau(r)=\frac{2rc_H-\ell^2+p^2}{2rh-r^2-\ell^2}
=\frac{48r-3}{4(20r-r^2-1)}\quad\text{at the benchmark}.
\end{equation}
For the derivative, write \(N_\tau=2rc_H-\ell^2+p^2\) and \(D_\tau=2rh-r^2-\ell^2>0\).
The quotient rule gives
\begin{align*}
N_\tau'D_\tau-N_\tau D_\tau'
&=2c_H(2rh-r^2-\ell^2)-(2rc_H-\ell^2+p^2)(2h-2r)\\
&=2c_Hr^2-2(\ell^2-p^2)r+2h(\ell^2-p^2)-2c_H\ell^2.
\end{align*}
Regroup the numerator as
\[
2c_H(r^2-\ell^2)+2(h-r)(\ell^2-p^2)>0.
\]
Both terms are positive because \(c_H>0\) and \(0<p<\ell<r<h\).
Thus \(\tau\) rises throughout this domain, even if the threshold lies outside \([0,1]\).
At the benchmark,
\begin{equation}\label{eq:03-taup}
\tau'(r)=\frac{3(8r^2-r+2)}{2(20r-r^2-1)^2}>0\qquad(1<r<10).
\end{equation}
At \(r=1\), \(\tau=5/8>1/2\).
Therefore an uninformative price never recruits a high-cost challenger along the benchmark path.

\begin{resultbox}[The two effects]
A stronger uniform incumbent raises \(\DT\) and lowers gross profit at every fixed belief.
It also raises \(\tau\) in this benchmark domain.
The top posterior must clear this higher threshold before a purchase recruits.
\end{resultbox}

\subsection{Benchmark numbers (Table 1 of the paper)}\label{sec:03_auction-table}

With \(h=10\), \(\ell=1\), \(p=1/2\) and \(c_H=6\), the formulas give Table~\ref{tab:03_auction-table1}.
Each entry is an exact fraction; the decimal is the paper's.

\begin{table}[!htbp]
\centering\small
\setlength{\tabcolsep}{4pt}
\begin{tabular}{@{}llccc@{}}
\toprule
Term & Formula & \(r_0=\reg{prim_r0}\) & \(r_1=\reg{prim_r1}\) & \(r_2=\reg{prim_r2}\)\\
\midrule
\(t_0\) & \(1/2-1/(4r)\) & \(7/24=\reg{p2_base_r0_t0}\) & \(5/12=\reg{p2_base_r1_t0}\) & \(31/72=\reg{p2_base_r2_t0}\)\\
\(t_H\) & \(r/2+1/(8r)\) & \(169/240=\reg{p2_base_r0_th}\) & \(37/24=\reg{p2_base_r1_th}\) & \(1321/720=\reg{p2_base_r2_th}\)\\
\(t_L\) & \(1-3/(8r)\) & \(11/16=\reg{p2_base_r0_tl}\) & \(7/8=\reg{p2_base_r1_tl}\) & \(43/48=\reg{p2_base_r2_tl}\)\\
\(\DT\) & \((r-1)^2/(2r)\) & \(1/60=\reg{Delta_r0}\) & \(2/3=\reg{Delta_r1}\) & \(169/180=\reg{Delta_r2}\)\\
\(g_H\) & \(10-t_H\) & \(2231/240=\reg{p2_base_r0_gh}\) & \(203/24=\reg{p2_base_r1_gh}\) & \(5879/720=\reg{p2_base_r2_gh}\)\\
\(g_L\) & \(3/(8r)\) & \(5/16=\reg{p2_base_r0_gl}\) & \(1/8=\reg{p2_base_r1_gl}\) & \(5/48=\reg{p2_base_r2_gl}\)\\
\(B_r(1/2)\) & \(5-r/4+1/(8r)\) & \(1153/240=\reg{Bprior_r0}\) & \(103/24=\reg{Bprior_r1_B}\) & \(2977/720=\reg{Bprior_r2}\)\\
\(\tau(r)\) & see the text & \(195/308=\reg{tau_r0}\) & \(141/200=\reg{tau_r1}\) & \(4245/5804=\reg{tau_r2}\)\\
\bottomrule
\end{tabular}
\caption{Acquisition payoffs in the benchmark \paper{Table~1}. They do not depend on the slot.}
\label{tab:03_auction-table1}
\end{table}

The row for \(B_r(1/2)\) uses \((g_H+g_L)/2=\tfrac12\bigl(10-\tfrac r2-\tfrac1{8r}+\tfrac3{8r}\bigr)
=5-\tfrac r4+\tfrac1{8r}\).

Note the size of the change from \(r_0\) to \(r_1\). The spread rises \(40\) times
(\((2/3)/(1/60)=40\)). Gross profit at the prior falls by \(41/80\), about \(10.7\) percent.
% verify/03_auction.py
The threshold rises from \reg{tau_r0} to \reg{tau_r1}, as \eqref{eq:03-taup} requires.

\pagebreak
\subsection{Proposition 1: the result for any incumbent distribution}\label{sec:03_auction-prop1}

\begin{resultbox}[Proposition 1 (competition and the two returns to information)]
Let \(R\) have a continuous distribution \(F\) on \([0,\bar r]\), with \(F(y)=1\) for \(y>\bar r\). For any
classes \(G_H\), \(G_L\) with finite means,
\[
\DT(F)=pF(p)D(p^-)+\int_p^\infty[1-F(y)]\,D(y)\,dy,\qquad
g_j(F)=\int_p^\infty[1-G_j(y)]\,F(y)\,dy. \tag*{\paper{(8)}}
\]
Let \(\tilde F\le F\) be a continuous first-order stochastic strengthening of \(F\).
\begin{enumerate}[label=(\alph*)]
\item For every pair of classes, the strengthening weakly lowers \(g_H\) and \(g_L\), and hence gross
profit at every fixed belief.
\item Under the class ranking \paper{(1)}, the strengthening weakly raises \(\DT\).
\end{enumerate}
The decrease in \(g_j\) is strict when \(\int_p^\infty[1-G_j(y)][F(y)-\tilde F(y)]\,dy>0\). The increase in
(b) is strict when \(\int_p^\infty[F(y)-\tilde F(y)]\,D(y)\,dy>0\).
\end{resultbox}

\paragraph{Step 1: write each payment as an integral of an indicator} \paper{(A.2)}. For \(\theta\ge p\),
\[
\max\{p,\min(R,\theta)\}=p+\int_p^\theta\1\{R>y\}\,dy,\qquad
(\theta-\max\{p,R\})_+=\int_p^\theta\1\{R<y\}\,dy.
\]
Check both identities case by case.
\begin{itemize}
\item \(R\le p\). Every \(y\in(p,\theta)\) exceeds \(R\). The first integral is \(0\), so the right side is
\(p\); the left side is \(\max\{p,\min(R,\theta)\}=p\). The second integral is \(\theta-p\); the left side
is \(\theta-\max\{p,R\}=\theta-p\).
\item \(p<R<\theta\). The first indicator is \(1\) on \((p,R)\), so the right side is \(p+(R-p)=R\); the
left side is \(\min(R,\theta)=R\). The second indicator is \(1\) on \((R,\theta)\), so the integral is
\(\theta-R\); the left side is \(\theta-R\).
\item \(R\ge\theta\). The first indicator is \(1\) on all of \((p,\theta)\), so the right side is
\(\theta\); the left side is \(\min(R,\theta)=\theta\). The second integral is \(0\); the left side is
\((\theta-R)_+=0\).
\end{itemize}
The two indicators add to one except at \(y=R\). Thus target proceeds plus challenger profit equal
\(p+(\theta-p)=\theta\) at every \(R\): \(t(\theta)+g(\theta)=\theta\) for \(\theta\ge p\).

\paragraph{Step 2: take the expectation over \(R\).} Each integrand is nonnegative and jointly measurable in
\((R,y)\). So Tonelli's theorem lets us move the expectation inside the integral. Since \(F\) is
continuous, \(\Pr(R=y)=0\), so \(\Pr(R>y)=1-F(y)\) and \(\Pr(R<y)=F(y)\). For \(\theta\ge p\):
\[
t(\theta)=p+\int_p^\theta[1-F(y)]\,dy,\qquad g(\theta)=\int_p^\theta F(y)\,dy.
\]
For \(\theta<p\): \(t(\theta)=t_0=p[1-F(p)]\) and \(g(\theta)=0\).

\paragraph{Step 3: one formula for every \(\theta\).} At \(\theta=p\), \(t\) jumps by \(pF(p)\).
With probability \(F(p)\), the incumbent fails the reserve.
An admissible challenger then buys at \(p\); an inadmissible one leaves the target unsold. Since \(t_0+pF(p)=p[1-F(p)]+pF(p)=p\),
\[
t(\theta)=t_0+pF(p)\,\1\{\theta\ge p\}+\int_p^\infty\1\{y\le\theta\}[1-F(y)]\,dy,\qquad
g(\theta)=\int_p^\infty\1\{y\le\theta\}F(y)\,dy.
\]

\paragraph{Step 4: take the expectation over the class.} Apply Tonelli's theorem again, with
\(\E_j\1\{\theta\ge p\}=1-G_j(p^-)\) and \(\Pr_j(\theta\ge y)=1-G_j(y^-)\). The functions \(G_j(y^-)\) and
\(G_j(y)\) differ only at the atoms of \(G_j\), which are countably many, a set of Lebesgue measure zero.
So \(G_j(y)\) can replace \(G_j(y^-)\) inside the integral:
\[
t_j=t_0+pF(p)\,[1-G_j(p^-)]+\int_p^\infty[1-G_j(y)][1-F(y)]\,dy,\qquad
g_j=\int_p^\infty[1-G_j(y)]\,F(y)\,dy.
\]
Each integral is finite, since \(0\le\int_p^\infty[1-G_j(y)]\,dy\le\E_j\theta<\infty\) and \(0\le F\le1\).

\paragraph{Step 5: the spread.} Subtract the class-\(L\) formula from the class-\(H\) formula. The term
\(t_0\) cancels. Since \([1-G_H]-[1-G_L]=G_L-G_H=D\),
\begin{align*}
\DT=t_H-t_L&=pF(p)\bigl[G_L(p^-)-G_H(p^-)\bigr]+\int_p^\infty D(y)[1-F(y)]\,dy\\
&=pF(p)D(p^-)+\int_p^\infty[1-F(y)]\,D(y)\,dy .
\end{align*}

\paragraph{Step 6: part (a).} In \(g_j\), the weight \(1-G_j\) on \(F\) is nonnegative. A pointwise fall from
\(F\) to \(\tilde F\) lowers \(g_j\) by
\[
g_j(F)-g_j(\tilde F)=\int_p^\infty[1-G_j(y)]\,[F(y)-\tilde F(y)]\,dy\ge0,
\]
strictly when this integral is positive. Gross profit at a fixed belief \(\mu\) is \((1-\mu)g_L+\mu g_H\),
an average with fixed weights, so it falls as well.

\paragraph{Step 7: part (b).} The difference of the spread formula at the two distributions is
\[
\DT(\tilde F)-\DT(F)=\int_p^\infty[F(y)-\tilde F(y)]\,D(y)\,dy-p\,[F(p)-\tilde F(p)]\,D(p^-).
\]
To see this, write \([1-\tilde F]-[1-F]=F-\tilde F\) inside the integral, and
\(p\tilde F(p)D(p^-)-pF(p)D(p^-)=-p[F(p)-\tilde F(p)]D(p^-)\). Under the class ranking, \(D(p^-)=0-0=0\),
so the second term is zero. Also \(D(y)\ge0\) and \(F(y)-\tilde F(y)\ge0\) for every \(y\), so the
integrand is nonnegative. The spread rises, strictly when the integral is positive. \hfill\(\square\)

\paragraph{Check with the benchmark.} For point masses with \(p<\ell\), \(G_L(y)=\1\{y\ge\ell\}\) and
\(G_H(y)=\1\{y\ge h\}\), so \(D(y)=\1\{\ell\le y<h\}\) and \(D(p^-)=0\). With \(F(y)=y/r\) on \([0,r]\)
and \(r<h\):
\begin{align*}
\DT&=\int_\ell^h[1-F(y)]\,dy=\int_\ell^r\Bigl(1-\frac yr\Bigr)dy=(r-\ell)-\frac{r^2-\ell^2}{2r}
=\frac{(r-\ell)(2r-r-\ell)}{2r}=\frac{(r-\ell)^2}{2r},\\
g_L&=\int_p^\ell\frac yr\,dy=\frac{\ell^2-p^2}{2r},\qquad
g_H=\int_p^r\frac yr\,dy+\int_r^h1\,dy=\frac{r^2-p^2}{2r}+h-r=h-\frac r2-\frac{p^2}{2r}.
\end{align*}
All three match \paper{(6)}. The script \texttt{verify/03\_auction.py} also checks \paper{(8)} for a
non-uniform \(F(y)=(y/\bar r)^2\) and classes with an atom: the formula and a direct double expectation
agree to twelve digits.

\begin{intuition}[slices of the value gap]
Read the integrand of \(\DT\) as slices \(D(y)\,dy\) of the class value gap at each level \(y>p\). The
slice at level \(y\) reaches the target when \(R>y\): the incumbent outbids the challenger at that level
and pays its bid, or the challenger pays at least \(y\). The slice stays with the challenger when \(R<y\).
A stronger incumbent raises \(\Pr(R>y)\) at every \(y\). So it moves slices from the challenger's profit
to target proceeds: the spread rises while gross profit falls.
\end{intuition}

\begin{caution}[scope of Proposition 1]
The kernel belongs to the cash second-price rule with reserve: an incumbent that wins pays the loser's bid.
Section~7.2 of the paper (\Cref{sec:11_welfare} here) changes the payment rule to Nash bargaining
with seller weight \(\omega\). The effect of competition on the spread then has the sign of \(1-2\omega\).
The zero-reserve second-price auction is the case \(\omega=0\).
\end{caution}

\subsection{Proposition 1(c): the weakest restriction}\label{sec:03_auction-prop1c}

\begin{resultbox}[Proposition 1(c)]
Let \(G_H\) and \(G_L\) have finite means. The spread weakly rises under every first-order stochastic
strengthening of every continuous \(F\) if and only if \(D(y)\ge0\) for every \(y>p\) and \(D(p^-)\le0\).
When neither class has an atom at \(p\), this condition is exactly first-order dominance above the reserve
together with \emph{equal admissibility}, \(G_H(p^-)=G_L(p^-)\).
\end{resultbox}

\paragraph{Step 1: the kernel.} Fix \(R\) and compute the expected class difference in proceeds.
\begin{itemize}
\item \(R<p\). The incumbent cannot bid. An admissible challenger buys at \(p\); an inadmissible one leaves
the target unsold. The target receives \(p\,\1\{\theta\ge p\}\), with mean \(p[1-G_j(p^-)]\) in class
\(j\). The class difference is \(p[G_L(p^-)-G_H(p^-)]=pD(p^-)\).
\item \(R\ge p\). The target receives \(\min\{R,\max\{\theta,p\}\}\). If \(\theta\ge p\), both bidders
bid at least \(p\), and the target gets \(\min(R,\theta)\). If \(\theta<p\), the incumbent buys alone at
\(p=\min\{R,p\}\). With \(Z=\max\{\theta,p\}\ge p\), the indicator form gives
\(\min\{R,Z\}=p+\int_p^R\1\{Z>y\}\,dy\). For \(y>p\), \(Z>y\) exactly when \(\theta>y\). So
\(\E_j[\min\{R,Z\}]=p+\int_p^R[1-G_j(y)]\,dy\), and the class difference is \(\int_p^RD(y)\,dy\).
\end{itemize}
So
\[
\xi(R)=\begin{cases} pD(p^-), & R<p,\\[2pt] \int_p^R D(y)\,dy, & R\ge p,\end{cases}
\qquad\text{and}\qquad \DT(F)=\E_F[\xi(R)].
\]

\paragraph{Step 2: when is \(\xi\) nondecreasing?} The kernel is bounded, since
\(\int_p^\infty|D(y)|\,dy\le\E_H\theta+\E_L\theta<\infty\). It is constant below \(p\). On \([p,\infty)\) it
is continuous, and its derivative is \(D\) almost everywhere. At \(p\) it moves from \(pD(p^-)\) to
\(\xi(p)=0\). So \(\xi\) is nondecreasing if and only if
\begin{itemize}
\item \(D\ge0\) almost everywhere above \(p\) (no fall on \([p,\infty)\)), and
\item \(pD(p^-)\le0\), that is \(D(p^-)\le0\) (no fall at the step).
\end{itemize}

\paragraph{Step 3: ``if''.} If \(\xi\) is nondecreasing, every first-order strengthening raises
\(\E_F[\xi(R)]\). One line: with the quantile coupling, \(\tilde R=\tilde F^{-1}(U)\ge F^{-1}(U)=R\) for a
uniform \(U\), so \(\xi(\tilde R)\ge\xi(R)\) at every \(U\).

\paragraph{Step 4: ``only if''.} Suppose \(\xi\) is not nondecreasing. Then there are values
\(0<R_1<R_2\) with \(\xi(R_1)>\xi(R_2)\). Continuity on each side of \(p\) lets us pick both different from
\(p\). Let \(F\) be uniform on \([R_1-\delta,R_1+\delta]\), and let \(\tilde F\) be the same law shifted
to \(R_2\). For \(0<\delta<\min\{R_1,(R_2-R_1)/2,|R_1-p|,|R_2-p|\}\), \(\tilde F\le F\) at every point: \(\tilde F>0\) only above
\(R_2-\delta\), where \(F=1\). As \(\delta\to0\), \(\DT(\tilde F)-\DT(F)\to\xi(R_2)-\xi(R_1)<0\). So some
strengthening lowers the spread. If the support must start at zero, mix the same small uniform component
on \([0,\bar r]\) into both laws. The mixtures stay continuous and ordered, and the difference keeps its
sign, scaled by a positive factor.

\paragraph{Step 5: the case without an atom at \(p\).} Distribution functions are right-continuous, so
\(D\) is right-continuous. If \(D(y_0)<0\) at some \(y_0>p\), then \(D<0\) on an interval
\([y_0,y_0+\varepsilon)\), a set of positive measure. So ``\(D\ge0\) almost everywhere above \(p\)'' is the
same as ``\(D(y)\ge0\) for every \(y>p\)'', and it gives \(D(p)=\lim_{y\downarrow p}D(y)\ge0\). Without
an atom at \(p\), \(D(p^-)=D(p)\ge0\). With \(D(p^-)\le0\), this forces \(D(p^-)=0\), which is equal
admissibility. \hfill\(\square\)

\paragraph{A worked example where the spread falls.} Take the benchmark classes with the exclusion reserve
\(p=\reg{rsv_points_excl}>\ell\). The low class never meets the reserve, so \(G_L(p^-)=1\),
\(G_H(p^-)=0\) and \(D(p^-)=1\). Then \(\xi(R)=p\) for \(R<p\): only class \(H\) buys, at \(p\). And
\(\xi(R)=R-p\) for \(p\le R<h\): class \(H\) wins and pays \(R\), while a class-\(L\) challenger cannot
bid, so the incumbent buys alone at \(p\). For the uniform incumbent on \([0,r]\) with \(p<r<h\),
\[
\DT(r)=\frac1r\Bigl[\int_0^p p\,dy+\int_p^r(y-p)\,dy\Bigr]=\frac{p^2+(r-p)^2/2}{r}.
\]
Its derivative is \(\bigl[(r-p)r-p^2-(r-p)^2/2\bigr]/r^2=(r^2-3p^2)/(2r^2)\), negative for
\(r<\sqrt3\,p\approx1.749\). % verify/03_auction.py
So from \(r=6/5\) to \(r=3/2\) the spread \emph{falls}, from \(6921/8000=0.865125\) to
\(7601/10000=0.7601\). % verify/03_auction.py
It rises again later and reaches \(20001/20000\) at \(r=3\). % verify/03_auction.py
The cause is the reserve term: a stronger incumbent fails the reserve less often, so the slice \(p\) that
only class \(H\) delivers below the reserve shrinks. Section~7.3 of the paper (\Cref{sec:11_welfare}
here) studies this reserve in full.

\subsection{The surplus split}\label{sec:03_auction-split}

Add the two lines of \paper{(8)}. First, \(g_H-g_L=\int_p^\infty\{[1-G_H]-[1-G_L]\}F\,dy
=\int_p^\infty D(y)F(y)\,dy\). So
\begin{align*}
\DT+g_H-g_L&=pF(p)D(p^-)+\int_p^\infty D(y)[1-F(y)]\,dy+\int_p^\infty D(y)F(y)\,dy\\
&=pF(p)D(p^-)+\int_p^\infty D(y)\,dy.
\end{align*}
Next, \(\max\{\theta,p\}=p+(\theta-p)_+\) and \(\E_j(\theta-p)_+=\int_p^\infty\Pr_j(\theta>y)\,dy
=\int_p^\infty[1-G_j(y)]\,dy\). So \(\int_p^\infty D(y)\,dy=\E_H[\max\{\theta,p\}]-\E_L[\max\{\theta,p\}]\),
and
\[
\DT(F)+g_H(F)-g_L(F)=pF(p)D(p^-)+\E_H[\max\{\theta,p\}]-\E_L[\max\{\theta,p\}].
\]
The reserve term vanishes under equal admissibility, and also when the incumbent always meets the reserve
(\(F(p)=0\)). Under the class ranking every value is at least \(p\), so \(\max\{\theta,p\}=\theta\) and
\[
\DT(F)+g_H(F)-g_L(F)=\E_H\theta-\E_L\theta\qquad\text{for every }F.
\]
In the benchmark the right side is \(h-\ell=9\) at every \(r\). Check at \(r_1=3\):
\(2/3+203/24-1/8=16/24+203/24-3/24=216/24=9\).

\begin{panelnote}[how to say it]
``A stronger incumbent does not change the value gap between the classes. It moves the gap from the
challenger's profit to target proceeds. The stock trades on the target's share of the gap, and the
challenger decides on its own share.''
\end{panelnote}

\subsection{Link to the old guide}\label{sec:03_auction-laplace}

The auction stage does not depend on the financial market.
Thus the six benchmark payoffs remain those of the old guide.
The new material is the formula for general classes, the reserve term, and Proposition~1(c).
\Cref{sec:10_robustness-sizes} covers the change in the financial market.
